\section{Discussion and Limitations}
\label{sec:discussion}

\noindent \textbf{Checking scope.}
\tool{} primarily targets requirements that can be mapped to protocol fields,
trigger conditions, and implementation behavior. It is well suited for checking
local conformance deviations in parsing, validation, negotiation, and error
handling. For global properties such as cross-session behavior or long-term
state evolution, \tool{} can only provide limited approximations and cannot
offer completeness or correctness guarantees; strict conclusions still require
separate modeling and verification.

\noindent \textbf{Natural-language ambiguity.}
Ambiguity in standard text can affect rule extraction and conformance
judgment. For example, a requirement such as "reject malformed messages" may
not clearly specify the relevant field combinations, compatibility exceptions,
or error-handling paths. In such cases, \tool{} may misclassify reasonable
compatibility behavior, or it may miss real defects. Therefore, for rules with
unclear semantic boundaries or rules that depend on implementation context, the
final conclusion still requires expert interpretation and maintainer feedback.
Our artifact preserves the source text, reasoning evidence, and uncertainty
notes to support subsequent review.

\noindent \textbf{LLM dependence.}
\tool{} uses LLMs for extraction, ranking, semantic filtering, reasoning, and
test planning. The framework constrains these calls with schemas, field spaces,
syntax evidence, and deterministic validators, but it does not eliminate model
error. We mitigate this risk by retaining uncertainty, using low-temperature
decoding, validating outputs, and requiring executable evidence for confirmed 
reports.

\noindent \textbf{Source availability.}
The current prototype relies on source code to locate variables, call
relationships, and validation logic, and therefore does not support
closed-source implementations. For implementations produced by code generators,
or implementations that encapsulate substantial parsing, validation, and
error-handling logic in macros, semantic cues in the source code may be unstable
or incomplete, which can affect the quality of specification-to-implementation
alignment.

\noindent \textbf{Language front-end.}
The current prototype mainly supports C/C++ protocol implementations. It
combines syntactic analysis with LLM-assisted filtering to extract code context
related to field accesses, conditional checks, and error handling from source
code. For other programming languages, the system still requires corresponding
language front-ends that provide equivalent code-analysis capabilities. In the
future, we plan to extend this capability by integrating front-end analysis
components for additional languages.

\noindent \textbf{Result adjudication.}
Protocol conformance judgments cannot always be directly determined by the
authors, because standards may be ambiguous and implementations may adopt
different behaviors for compatibility. Therefore, system outputs should not be
treated as final conclusions by themselves. We adjudicate reports as valid
issues, false positives, or pending based on maintainer feedback, independent
review, and reproducible tests. This also means that some results depend on
external feedback and follow-up validation, and the current statistics only
reflect the adjudication status at the time of evaluation.

\noindent \textbf{Cost and Human Effort.}
\tool{} reduces model cost by caching standard-extraction results,
compressing code contexts, and invoking the validation stage only for
suspicious samples or samples with insufficient evidence. Nevertheless,
multi-round reasoning and validation still incur higher model-calling costs
than one-shot prompting, and report confirmation still requires some manual
involvement. The manual effort mainly comes from reviewing candidate
inconsistencies, including checking report accuracy, filtering a small number
of false positives, confirming the execution results, and submitting reports to
maintainers with follow-up communication.
% Deployment in large protocol ecosystems still requires
% prioritization; one practical strategy is to first analyze check targets with
% high security relevance, dense historical defects, or strong rule constraints.

\noindent \textbf{Threats to validity.}
Our evaluation may be affected by manual adjudication and model stability.
First, manually labeled ground truth may be influenced by reviewers'
interpretations of ambiguous standards. Second, LLM nondeterminism may change
extraction or reasoning results across runs. We mitigate these threats by using
two-author adjudication, preserving raw evidence, separating pending reports
from confirmed fixes, and measuring variance across repeated GPT-5.4 runs and
different model families.








% \textbf{What would make a report actionable?}
% Maintainers are unlikely to act on a vague statement that code and standard
% ``seem inconsistent.'' A useful report should include the exact standard span,
% the normalized rule, the relevant field and implementation variable, the code
% path that accepts or rejects the input, the build configuration, and a minimal
% reproduction script. \tool{} emits this bundle and stores it under a content
% hash so that reviewers can inspect the evidence without rerunning the full
% pipeline.


% \textbf{Generalization beyond protocols.}
% The field-centric idea may transfer to other domains where natural-language
% documents define structured objects, such as file formats, API contracts, or
% hardware specifications. The representation would need different object spaces:
% API parameters instead of protocol fields, register bits instead of message
% fields, or configuration keys instead of transport parameters. This paper
% focuses on network protocols because their standards and implementations make
% the alignment problem both important and concrete.
